% Complete Chapter 7 introduction, original PDF 91-92 / printed 81-82.
% block: finite-element-forms-chapter-heading page:91 type:chapter
\clearpage\setcounter{section}{6}
\section{{Finite element differential forms}}\label{chap:finite-element-differential-forms}

% block: finite-element-forms-hodge-problem page:91 type:ordinary-paragraph
{Consider now a Lipschitz polyhedral domain $\Omega\subset\mathbb R^n$ and its $L^2$ de Rham complex \eqref{eq:l2-manifold-de-rham-complex} with domain complex \eqref{eq:manifold-de-rham-domain-complex}. For any $k$, the segment}
\begin{equation}\label{eq:finite-element-forms-short-complex}
H\Lambda^{k-1}(\Omega)\xrightarrow{d}H\Lambda^k(\Omega)\xrightarrow{d}H\Lambda^{k+1}(\Omega)
\end{equation}
{determines a Hodge Laplace boundary value problem involving the two exterior derivatives and their adjoints, as discussed in Section~\ref{sec:abstract-hodge-laplacian}. In Theorem~\ref{thm:exterior-codifferential-adjoints} we have identified the adjoint of $(d,H\Lambda^k)$ as $(\delta,\mathring H^*\Lambda^k)$, so the strong form of the $k$-form Hodge Laplace boundary value problem is}
\begin{equation}\label{eq:finite-element-forms-hodge-strong}
\begin{gathered}
(\delta d+d\delta)u=f-P_{\mathsf H}^kf\text{ in }\Omega,\\
\operatorname{tr}\star u=0,\ \operatorname{tr}\star du=0\text{ on }\partial\Omega,\qquad u\perp\mathsf H^k,
\end{gathered}
\end{equation}
{where $\mathsf H^k$ are the harmonic $k$-forms on $\Omega$, determined by the differential equations $dp=0$, $\delta p=0$ and the boundary condition $\operatorname{tr}\star p=0$. In the case of three dimensions, we have interpreted this problem, and related ones, in Section~\ref{sec:hodge-r3}. In Chapter~\ref{chap:approximation-hilbert-complexes} we derived the requirements needed for a successful discretization of this problem using the mixed formulation and Galerkin's method. For this we need to construct Galerkin subspaces $V_h^{k-1}$ of $H\Lambda^{k-1}(\Omega)$ and $V_h^k$ of $H\Lambda^k(\Omega)$. The key requirements of these subspaces, besides from being efficiently implementable, are that they afford good approximation, form a subcomplex, and admit a bounded cochain projection. In that case, as shown in Chapter~\ref{chap:approximation-hilbert-complexes}, the mixed Galerkin method is consistent and stable and we obtain convergence with an optimal rate.}

% block: finite-element-forms-construction-goal page:91 type:ordinary-paragraph
{In this chapter, we shall obtain efficiently implementable Galerkin subspaces of $H\Lambda^k$ by constructing them as spaces of finite elements. The desired degree of approximation will then be obtained by using a sufficiently fine mesh and/or sufficiently high degree piecewise polynomial shape functions. The challenge is to ensure that the spaces form a subcomplex admitting a bounded cochain projection.}

% block: complete-trimmed-families pages:91,92 type:ordinary-paragraph
{Specifically, given any simplicial triangulation $\mathcal T_h$ of $\Omega$ (in any number of dimensions $n\geq1$), any form degree $0\leq k\leq n$, and any polynomial degree $r\geq1$, we shall construct two finite element subspaces of $H\Lambda^k(\Omega)$, denoted}
\[
\mathcal P_r\Lambda^k(\mathcal T_h)\quad{\text{and}}\quad\mathcal P_r^-\Lambda^k(\mathcal T_h).
\]
{The former space uses as shape functions the \emph{complete} space of polynomial $k$-forms with degree at most $r$, while the latter uses a somewhat smaller space of polynomials which, following \BookCite{29}, we refer to as \emph{trimmed}. Then we show several ways to select two such spaces, one for $(k-1)$-forms, the other for $k$-forms, to solve the Hodge Laplacian. These complete and trimmed families of finite element spaces of differential forms can be viewed as the most canonical for discretization of the de Rham complex. This is one of the important conclusions of FEEC.}

% block: finite-element-forms-chapter-outline page:92 type:ordinary-paragraph
{The chapter is organized as follows. Before defining the complete and trimmed spaces $\mathcal P_r\Lambda^k(\mathcal T_h)$ and $\mathcal P_r^-\Lambda^k(\mathcal T_h)$ in Section~\ref{sec:finite-element-differential-forms}, we establish some properties of differential forms with polynomial coefficients in Sections~\ref{sec:complete-polynomial-forms}--\ref{sec:trimmed-polynomial-forms}. Then, in Section~\ref{sec:finite-element-space-properties}, we show how these spaces may be combined into complexes which fulfill the hypotheses required for the error analysis in Chapter~\ref{chap:approximation-hilbert-complexes}, thus establishing a full set of error bounds for our methods. In Section~\ref{sec:cubical-elements} we briefly describe analogous results for finite elements based on meshes of boxes rather than simplices. Finally, in Section~\ref{sec:historical-antecedents}, we relate the elements described in this chapter to various finite elements in the literature.}

% block: finite-element-forms-notation page:92 type:ordinary-paragraph
{Before proceeding, we make a brief remark on notation. The finite element spaces $\mathcal P_r\Lambda^k(\mathcal T_h)$ and $\mathcal P_r^-\Lambda^k(\mathcal T_h)$ which we will define in the next section consist of piecewise polynomial differential $k$-forms defined over the domain $\Omega$. We will give a unified construction of the spaces, but each space depends on the polynomial degree $r\geq0$, the form degree $k$ between $0$ and $n=\dim\Omega$, and the simplicial mesh $\mathcal T_h$. We shall also use the notation $\mathcal P_r\Lambda^k$ and $\mathcal P_r^-\Lambda^k$ for the spaces of complete and trimmed polynomial (not piecewise polynomial) differential $k$-forms of degree $r$ that are used on each simplex to construct the finite element spaces. These spaces of polynomials are defined for each $n$ and for each $r\geq0$ and $0\leq k\leq n$. Finally, we write $\mathcal P_r\Lambda^k(T)$ and $\mathcal P_r^-\Lambda^k(T)$ for the spaces of restrictions of these polynomials to a single simplex $T$. Thus our notation is analogous to a common notation for standard Lagrange finite elements of degree $r$: $\mathcal P_r(\mathcal T_h)$ for the finite element space, $\mathcal P_r$ for the space of polynomials of degree $r$, and $\mathcal P_r(T)$ for their restriction to an element.}
